\documentclass[11pt,a4paper]{article}

% Packages
\usepackage[utf8]{inputenc}
\usepackage[T1]{fontenc}
\usepackage{amsmath,amssymb,amsthm}
\usepackage{mathtools}
\usepackage{graphicx}
\usepackage{hyperref}
\usepackage{cleveref}
\usepackage{booktabs}
\usepackage{array}
\usepackage{xcolor}
\usepackage{microtype}
\usepackage[numbers]{natbib}

% Page layout
\usepackage[margin=1in]{geometry}
\emergencystretch=1em
\tolerance=1000

% Theorem environments
\theoremstyle{plain}
\newtheorem{theorem}{Theorem}[section]
\newtheorem{lemma}[theorem]{Lemma}
\newtheorem{corollary}[theorem]{Corollary}
\newtheorem{proposition}[theorem]{Proposition}
\newtheorem{conjecture}[theorem]{Conjecture}

\theoremstyle{definition}
\newtheorem{definition}[theorem]{Definition}
\newtheorem{remark}[theorem]{Remark}

% Notation shortcuts
\newcommand{\F}{\mathbb{F}}
\newcommand{\LL}{\mathcal{L}}
\newcommand{\E}{\mathbb{E}}
\newcommand{\Prb}{\Pr}
\DeclareMathOperator{\Lagr}{Lagr}
\DeclareMathOperator{\SDA}{SDA}
\DeclareMathOperator{\VSTAT}{VSTAT}
\DeclareMathOperator{\Adv}{Adv}
\DeclareMathOperator{\Var}{Var}
\newcommand{\etal}{\textit{et al.}}
\newcommand{\R}{\mathbb{R}}
\newcommand{\LSN}{\mathsf{LSN}}

\title{Lagrangian Subspace Noise: \\ Cryptographic Constructions}
\author{Kwanghoo Choo\thanks{ORCID: \href{https://orcid.org/0009-0005-5682-8098}{0009-0005-5682-8098}. Email: \texttt{gattochoo@gmail.com}. \textcopyright~2026 Kwanghoo Choo. Licensed under CC BY 4.0.}\\ \normalsize Independent Researcher}
\date{June 2026}

\begin{document}

\maketitle

\begin{abstract}
This is the constructions companion to the core paper \emph{Lagrangian Subspace Noise: Exact Statistical-Query Lower Bounds and the Linear-Reduction Landscape} (same author; this repository), where the hardness analysis of the Lagrangian Subspace Noise (LSN) problem --- exact statistical-query lower bounds, the near-complete linear-reduction landscape, the information-theoretic separations, and the security parameters --- is developed. LSN is a variant of Learning Parity with Noise whose secret is a Lagrangian subspace of a symplectic space over $\F_2$. Building on that hardness analysis, we develop a post-quantum zero-knowledge construction: \textbf{LSN-SNARK}, a preprocessing SNARK for the membership relation $x \in L$ whose membership circuit is exactly $n^2$ native-$\F_2$ AND-gate constraints ($4{,}225$ at $n=65$; $O(n^2)$ R1CS overall, far smaller than ZK-lattice verification circuits) via the Siegel-chart (graph-Lagrangian) representation. A key-encapsulation mechanism is more delicate: the natural reconciliation construction masks the codeword with publicly computable data (hence is insecure), and the constant-noise regime blocks the standard sparse-secret dual-LPN route, so we leave a sound LSN-KEM to future work. All security claims are classified as theorem, evidence, or conjecture; the construction inherits LSN's open-assumption status as analysed in the core paper.
\end{abstract}

\paragraph{\textbf{Keywords.}} Post-quantum cryptography, Learning Parity with Noise, symplectic geometry, zero-knowledge proofs, SNARKs.

\section{Introduction}
\label{sec:intro}

\subsection{Background}
\label{subsec:background}

The Learning Parity with Noise (LPN) problem, introduced by Blum \etal\ \cite{BFKL94}, asks to recover a secret linear function $s \in \F_2^k$ given noisy inner products $(a_i, \langle a_i, s \rangle \oplus e_i)$. Despite decades of cryptanalytic effort, no polynomial-time algorithm is known for LPN with constant noise rate. The best classical attacks run in sub-exponential time $2^{O(k / \log k)}$ using the Blum--Kalai--Wasserman (BKW) algorithm \cite{BKW03} and its modern refinements \cite{LF06,Lyu05,EKM17,GJL14,Kir11}. Grover's algorithm reduces the naive key-search cost from $2^k$ to $2^{k/2}$ \cite{Gro96}, and quantum variants of BKW remain the dominant quantum threat; no exponential quantum speed-up over the best classical methods is known.

LPN has been the foundation for a wide range of cryptographic primitives, including OWFs, PKE, OT, and MPC \cite{ABW10,Ale11,DGM12,BCGI18}. Its main theoretical limitation is the absence of a worst-case hardness reduction comparable to Regev's LWE reduction from lattice problems \cite{Reg05}. LPN hardness is therefore treated as a \emph{primitive assumption}, accepted on the basis of sustained cryptanalytic resistance rather than reducibility.

Post-quantum cryptography currently rests on a handful of established hardness families: lattices (LWE/SIS), error-correcting codes (LPN/syndrome decoding), multivariate quadratic equations (MQ), hash functions, and isogenies, with tensor and group-action problems emerging as further candidates. Shor's algorithm \cite{Sho94} breaks the integer-factoring and discrete-logarithm assumptions underlying classical public-key cryptography, motivating the search for quantum-resistant alternatives. A genuinely new family --- structurally distinct and not reducible to the existing ones --- has been sought since the inception of the field. LSN is a recent candidate for such a family, studied here without asserting its standing among the established ones.

\subsection{The LSN Problem}
\label{subsec:lsn}

We propose \emph{Lagrangian Subspace Noise} (LSN); its name and motivation come from the \emph{Learning Stabilizers with Noise} problem of \cite{PQS26} and the \emph{symplectic LPN} problem of \cite{KLP+25,LPQR26}, and the precise relation between our membership formulation and theirs is discussed in \Cref{subsec:two-forms}. The secret is not a linear function but a Lagrangian subspace $L \subset \F_2^{2n}$ of a symplectic vector space. The learner receives samples
\[
(a_i, b_i) \quad \text{where} \quad b_i = \mathbf{1}_L(a_i) \oplus e_i,
\]
with $a_i \sim \F_2^{2n}$ uniform, $e_i \sim \operatorname{Bernoulli}(p)$, and $\mathbf{1}_L$ the indicator of $L$. The secret space is the Lagrangian Grassmannian:
\begin{equation}
\label{eq:lagr-count}
|\Lagr(2n, \F_2)| = \prod_{i=1}^n (2^i + 1) = 2^{n(n+1)/2 + O(1)}.
\end{equation}
This is exponentially larger than LPN's $2^n$ secrets, yet it is not an unstructured space: every secret satisfies the strong isotropy constraint $L = L^{\perp_\Omega}$. It is exactly this tension --- exponential size versus rigid structure --- that creates both the hardness and the technical difficulty of LSN.

Recent independent work has put the stabilizer-decoding form of the problem on a firm footing: Poremba--Quek--Shor \cite{PQS26} introduced the Learning-Stabilizers-with-Noise problem (which contains LPN as a special case) and proved a worst-case-to-average-case reduction for a restricted (non-standard) variant of the problem; Khesin \etal\ \cite{KLP+25} proved that constant-rate LPN reduces to \emph{their} stabilizer-decoding LSN even for a single logical qubit; and Lu \etal\ \cite{LPQR26} established strong barriers against reducing symplectic LPN to standard LPN (for linear reductions, in the constant-rate sample regime $m=\Theta(n)$; the core paper develops a complementary near-complete landscape). The precise relation between their LSN and our membership formulation is stated as an open positioning item in \Cref{subsec:two-forms}.

\subsection{Scope and contributions}
\label{subsec:contributions}

This is the constructions companion. The hardness analysis the constructions rely on --- the exact pairwise-correlation formula, the unconditional $\Omega(2^n)$ and conditional $2^{2n-O(1)}$ statistical-query lower bounds, the near-complete linear-reduction landscape, the information-theoretic separation of the membership, batch, and sympLPN formulations, and the derivation of the security parameters --- is developed in the core paper \emph{Lagrangian Subspace Noise: Exact Statistical-Query Lower Bounds and the Linear-Reduction Landscape} (same author; this repository), to which we refer for all such results. Building on it, this report contributes:

\begin{enumerate}
    \item \textbf{LSN-SNARK (\Cref{sec:snark}).} A preprocessing SNARK for the LSN membership relation $x \in L$ whose membership circuit is exactly $n^2$ native-$\F_2$ AND-gates ($O(n^2)$ R1CS overall) via the Siegel-chart (graph-Lagrangian) representation. Verified counts: $64$ for $n=8$, $1{,}681$ for $n=41$, $4{,}225$ for $n=65$ (and $n$ constraints --- $8$, $41$, $65$ --- with public $a$).
    \item \textbf{On LSN-KEM (\Cref{sec:kem}).} We document why a KEM from LSN is delicate: the natural reconciliation construction masks the encapsulated codeword with publicly computable data (hence offers no security), and at constant noise $p=\tfrac14$ the standard sparse-secret dual-LPN route hits a capacity barrier. A sound construction and proof are left to future work.
\end{enumerate}

Every claim is classified as theorem, evidence, or conjecture. The constructions inherit LSN's open-assumption status; the standard-model gap (the absence of a non-vacuous LPN(low-noise)$\to$LSN or LWE$\to$LSN reduction) is discussed in the core paper.

\section{Preliminaries}
\label{sec:prelim}

\subsection{Symplectic Vector Spaces}
\label{subsec:symplectic}

Let $V = \F_2^{2n}$ with standard symplectic form
\[
\Omega(x,y) = \sum_{i=1}^n (x_i y_{n+i} + x_{n+i} y_i),
\]
\[
\text{i.e. } \Omega(x,y) = x^{\!\top} J y \text{ with } J = \begin{psmallmatrix} 0 & I_n \\ I_n & 0 \end{psmallmatrix} \text{ (char 2)}.
\]
A subspace $L \subset V$ is \emph{isotropic} if $\Omega|_{L \times L} = 0$, and \emph{Lagrangian} if it is isotropic of dimension $n$; equivalently $L = L^{\perp_\Omega}$. We write $\Lagr(2n, \F_2)$ for the Lagrangian Grassmannian.

\begin{proposition}[Lagrangian count \cite{AG04}]
$|\Lagr(2n, \F_2)| = \prod_{i=1}^n (2^i + 1) = 2^{n(n+1)/2 + O(1)}$.
\end{proposition}

For $L, L' \in \Lagr(2n)$ let $j = \dim(L \cap L')$ denote the intersection dimension; every value $0 \le j \le n$ occurs. The distribution of $j$ is governed by the $q$-binomial coefficients at $q=2$:
\begin{equation}
\label{eq:j-distribution}
\Pr[j=k] = \frac{{n \brack k}_2 \cdot 2^{(n-k)(n-k+1)/2}}{|\Lagr(2n,\F_2)|}.
\end{equation}

\subsection{Fourier Transform on the Symplectic Space}
\label{subsec:fourier}

The unnormalized $\Omega$-Fourier transform of $f \colon \F_2^{2n} \to \R$ is
\[
\mathcal{F}_\Omega[f](\xi) = \sum_{x \in \F_2^{2n}} f(x) (-1)^{\Omega(x, \xi)}.
\]

\begin{lemma}[Self-duality]
For every $L \in \Lagr(2n,\F_2)$,
\[
\mathcal{F}_\Omega[\mathbf{1}_L] = 2^n \cdot \mathbf{1}_L.
\]
\end{lemma}

\begin{proof}
If $\xi \in L$, then $\Omega(x,\xi)=0$ for all $x \in L$, so the sum is $|L|=2^n$. If $\xi \notin L$, choose $x_0 \in L$ with $\Omega(x_0,\xi)=1$; the involution $x \mapsto x+x_0$ pairs terms of opposite sign, so the sum is $0$.
\end{proof}

\subsection{Statistical Query Model}
\label{subsec:sq}

Let $\mathcal{D}$ be a class of distributions and $D_0$ a reference distribution. The \emph{statistical dimension with average correlation} \cite{FGR+17} is
\[
\SDA(B(\mathcal{D},D_0),\gamma) = \max\Bigl\{d : \forall S \subseteq \mathcal{D}, |S| \geq \frac{|\mathcal{D}|}{d} \Rightarrow \frac{1}{|S|^2}\sum_{D,D' \in S} |\langle D,D' \rangle| \leq \gamma\Bigr\}.
\]

\begin{theorem}[Feldman \etal\ \cite{FGR+17}]
\label{thm:feldman}
If $\SDA(B(\mathcal{D},D_0),\gamma) = d$, any SQ algorithm distinguishing $\mathcal{D}$ from $D_0$ with success probability $\alpha > 1/2$ requires $q \geq (2\alpha - 1)d$ queries to $\VSTAT(1/(3\gamma))$.
\end{theorem}

The theorem applies to randomized \emph{adaptive} SQ algorithms, so our lower bound will inherit adaptivity automatically.

\subsection{Notation}
\label{subsec:notation}

We write $D_L$ for the LSN distribution with secret $L$ and noise rate $p$, and $D_0$ for the noise-only distribution in which $b \sim \operatorname{Bernoulli}(p)$ independently of $a$. All logarithms are base $2$ unless stated otherwise.


\section{Related Work}
\label{sec:related}

\paragraph{LPN foundations and cryptanalysis.}
LPN was introduced by Blum, Furst, Kearns, and Lipton \cite{BFKL94} as a foundation for cryptographic primitives. The dominant classical attack is BKW \cite{BKW03}, later optimized with fast Walsh--Hadamard transforms \cite{LF06}, covering codes \cite{GJL14}, and hybrid information-set-decoding techniques \cite{EKM17}. Lyubashevsky \cite{Lyu05} gave a polynomial-sample variant, and a comprehensive unified attack framework was developed by Esser, K{\"u}bler, and May \cite{EKM17}. Quantum speed-ups are limited to Grover \cite{Gro96} and quantum variants of BKW; no exponential improvement over the best classical methods is known.

\paragraph{LPN variants.}
Many structured variants have been proposed to improve efficiency: Ring-LPN \cite{HKL+12}, LPN with structured (regular) noise \cite{BCGI18}, and the Learning with Rounding (LWR) problem \cite{BPR12}. All of these retain a secret space of size $2^{O(n)}$. LSN differs fundamentally in that its secret space has size $2^{n^2/2 + O(n)}$ while remaining polynomial-time sampleable via symplectic linear algebra.

\paragraph{Symplectic and quantum coding.}
Stabilizer codes, the quantum analogue of classical linear codes, are naturally described by Lagrangian subspaces of $\F_2^{2n}$ \cite{Got97,CRSS98,NC00}. Quantum state tomography of stabilizer states under depolarizing noise is a closely related learning problem \cite{GL22}. The recent independent work of Khesin \etal\ \cite{KLP+25} proved that constant-rate LPN reduces to \emph{their} stabilizer-decoding LSN even for $k=1$ (the relation to our membership formulation is an open positioning item; \Cref{subsec:two-forms}), while Lu \etal\ \cite{LPQR26} proved that linear reductions cannot reduce sympLPN to LPN in the constant-rate sample regime $m=\Theta(n)$; for $m=\omega(n)$ their bound shows error weight $>1/2-\delta$, which they note is not sufficient to fully rule out decodability (though they expect the barrier to extend to all $m=\operatorname{poly}(n)$).

\paragraph{Statistical query hardness.}
The SQ model was introduced by Kearns \cite{Kea98} and developed by Blum \etal\ \cite{BFJ+94}. Feldman, Grigorescu, Reyzin, and Vempala \cite{FGR+17} gave the general SDA framework we use, and Diakonikolas \etal\ \cite{DKS17} applied it to high-dimensional robust estimation. The known quantum approaches over classical samples (Grover search and quantum BKW) give no exponential speed-up over the best classical methods, so they do not affect the asymptotic security parameter. We do not, however, prove a quantum lower bound: the SQ bound is evidence, not proof, of quantum hardness (the core paper treats quantum hardness as a conjecture).

\paragraph{Code-based cryptography.}
Code-based cryptography dates to McEliece \cite{McE78} and Prange \cite{Pra62}; modern code-based KEM candidates include Classic McEliece \cite{ABD+21}, HQC \cite{AAC+22}, and BIKE, whose standardization status is discussed below.

\paragraph{Post-quantum KEMs.}
The NIST PQC standardization process \cite{NIS22} produced ML-KEM (FIPS~203, standardizing CRYSTALS-Kyber \cite{SAB+22}) and selected HQC \cite{AAC+22} for additional standardization; Classic McEliece \cite{ABD+21} remains a code-based candidate but is not a FIPS KEM standard. LSN targets a fundamentally different assumption from all of these.

\section{The LSN Problem}
\label{sec:lsn-problem}

\begin{definition}[Membership-LSN$_{n,m,p}$]
\label{def:lsn}
Let $L \subset \F_2^{2n}$ be uniformly random Lagrangian. The LSN distribution $D_L$ over $\F_2^{2n} \times \F_2$ is generated by sampling $a \sim \F_2^{2n}$, $e \sim \operatorname{Bernoulli}(p)$, and outputting $(a, b)$ where $b = \mathbf{1}_L(a) \oplus e$. Given $m$ independent samples from $D_L$, recover $L$.
\end{definition}

\paragraph{Search vs.~decision.} All hardness results in this paper are stated for the \emph{decisional} problem: distinguishing $D_L$ from the noise-only distribution $D_0$ (where $b \sim \operatorname{Bernoulli}(p)$ independently of $a$). Search-to-decision for membership-LSN does \emph{not} follow automatically from the LPN template: the secret is a Lagrangian subspace, not a linear function, and verifying a candidate $L'$ against a decision oracle only checks consistency --- it does not by itself search the $2^{\Theta(n^2)}$-size secret space. We do not prove a search-to-decision reduction here, and treat decisional LSN as the primitive assumption.

\subsection{Two Formulations}
\label{subsec:two-forms}

Definition~\ref{def:lsn} is the \emph{membership} (oracle) formulation: the query point $a$ is freshly random for every sample, and the adversary sees it only together with its label. This is the natural form for statistical-query and information-theoretic analysis.

Cryptographic constructions and external hardness results use the \emph{batch} formulation, where all query points are fixed in advance.

\begin{definition}[Batch-LSN$_{n,m,p}$]
\label{def:lsn-batch}
Let $L \subset \F_2^{2n}$ be uniformly random Lagrangian, and let $A = (a_1,\dots,a_m) \in (\F_2^{2n})^m$ be a fixed sequence of query points (public parameters). The batch distribution $D_L^{(A)}$ outputs labels $b = (b_1,\dots,b_m)$ where $b_j = \mathbf{1}_L(a_j) \oplus e_j$ with independent $e_j \sim \operatorname{Bernoulli}(p)$. Given $A$ and $b$, recover $L$.
\end{definition}

The columns of $A$ may be uniform random, pseudorandom (e.g.\ seed-derived in an application), or structured.

\begin{definition}[sympLPN$_{n,p}$ \cite{KLP+25,LPQR26}]
\label{def:symplpn}
Let $A \in \F_2^{2n \times n}$ be a public matrix whose columns are \emph{isotropic}: $S_A := A^{\top}\Omega A = 0$ (every pair of columns is orthogonal under the symplectic form). Let $x \in \F_2^n$ be a secret vector, $e \sim \operatorname{Bernoulli}(p)^{2n}$ a noise vector, and $y = Ax + e \in \F_2^{2n}$. Given $(A,y)$, recover $x$.
This is the special case $k=n$ of LPQR26's general sympLPN$(k,n,p)$; following LPQR26, the error is drawn from the symplectic representation of the depolarizing distribution, and for simplicity of analysis we state the problem with independent Bernoulli noise, which preserves all barrier results.
\end{definition}

\paragraph{Relationship.}
\begin{enumerate}
\item \textbf{Batch-LSN with uniform $A$ $\equiv$ membership-LSN.} When the columns of $A$ are uniform and independent, the adversary's view $(A,b)$ in batch-LSN is identical in distribution to $m$ independent samples from the membership oracle. Hence SQ lower bounds and information-theoretic floors proved for membership-LSN transfer \emph{verbatim} to batch-LSN with uniform~$A$.
\item \textbf{Pseudorandom $A$.} An application may derive the public sample points from a short seed via a random oracle. In the random-oracle model the derived points are uniform, so it reduces directly to batch-LSN with uniform $A$ (no separate pseudorandomness assumption).
\item \textbf{sympLPN is a different problem.} True sympLPN (\Cref{def:symplpn}) has linear labels $y = Ax + e$ and secret $x$; the isotropic constraint $S_A = 0$ is a public quadratic relation on the public parameters. It is not equivalent to batch-LSN (membership labels) without a pinned bridge theorem. LPQR26 study sympLPN; the transport theorems of the core paper use only the $S_A = 0$ structure on the public matrix and therefore apply to any formulation with isotropic columns, regardless of the label structure.
\item \textbf{Open positioning item: our membership-LSN vs.\ \linebreak[0] the LSN of \cite{KLP+25,LPQR26}.} Their (classical) LSN has a \emph{public} matrix $[A \mid B]$ --- the $n$ columns of $A$ are pairwise symplectically orthogonal and span a Lagrangian, and the $k$ columns of $B$ are likewise mutually isotropic and jointly independent of $A$; the concatenated matrix is \emph{not} itself isotropic, since $n$ is the maximal isotropic dimension --- a junk register $x$, and a secret logical string $y$. Each sample is a noisy codeword $[A \mid B] \cdot [x; y] + e$ under symplectic depolarizing noise, with a fresh public matrix per sample in the multi-sample variant $\mathrm{LSN}^m$. Our membership-LSN has no public matrix and the secret is the Lagrangian itself. The two formulations are not known to be equivalent, and we do not claim the external hardness results (e.g.\ constant-rate LPN $\le$ their LSN) for our membership formulation. Establishing or refuting this bridge is an open precision item (discussed in the core paper). We note for context that the core paper \emph{does} prove the cognate \emph{sympLPN} formulation (vector secret $x\in\F_2^n$, isotropic public matrix) inherits a constant-rate-LPN hardness floor, and that its symplectic structure is cryptanalytically inert; this anchors the hardness of the LSN \emph{family} in LPN, but it concerns a different formulation from the membership/batch-LSN our constructions use and does not transfer to them directly.
\end{enumerate}

Throughout the paper we state which formulation each theorem addresses; unless noted otherwise, ``LSN'' refers to the membership form.

\subsection{Information-Theoretic Floors}
\label{sec:info-floors}

The per-sample information content of membership-LSN is exponentially small, which makes the membership form statistically trivial at polynomial sample complexity and \emph{re-points} the hardness question to the batch formulation.

\begin{proposition}[Per-sample mutual information]
\label{prop:per-sample-mi}
For membership-LSN with noise rate $p$,
\[
I\bigl(L; (a,b)\bigr) \;=\; \Theta(2^{-n}).
\]
\end{proposition}

\begin{proof}
Given $L$, the label is $b = \mathbf{1}_L(a) \oplus e$ with $e \sim \operatorname{Bernoulli}(p)$, so $H(b \mid a, L) = H_2(p)$ exactly.  For $a \neq 0$, a uniformly random $a$ lies in $L$ with probability $(2^n-1)/(2^{2n}-1) = 2^{-n} + O(2^{-2n})$, hence $\Pr[b=1 \mid a] = p + 2^{-n}(1-2p) + O(2^{-2n})$; for $a=0$ we have $\Pr[b=1 \mid a=0] = 1-p$, and its entropy is $H_2(1-p) = H_2(p)$, but this occurs with probability $2^{-2n}$ and contributes negligibly.  Thus $H(b \mid a) = H_2(p + 2^{-n}(1-2p)) + O(2^{-2n})$ and the mutual information is $\Theta(2^{-n})$ by the Taylor expansion $H_2(p+\varepsilon) = H_2(p) + \varepsilon\log_2\frac{1-p}{p} + O(\varepsilon^2)$.
\end{proof}

\begin{proposition}[$\chi^2$ divergence and sample complexity]
\label{prop:chi2-sample}
For any fixed Lagrangian $L$, the $\chi^2$ divergence from $D_0$ is
\[
\chi^2(D_L \| D_0) \;=\; \frac{(1-2p)^2}{p(1-p)} \cdot 2^{-n}.
\]
Consequently any distinguisher (even computationally unbounded) needs $m = \Omega(2^n)$ samples to tell $D_L$ from $D_0$ with constant advantage, and any search algorithm needs $m = \Omega(n^2 \cdot 2^n)$ samples to recover $L$ with constant probability.
\end{proposition}

\begin{proof}
From a direct likelihood-ratio calculation (given in the core paper), the squared likelihood ratio has expectation $\E_{D_0}[\ell_L^2] = \frac{(1-2p)^2}{p(1-p)} \cdot 2^{-n}$.  Since $\E_{D_0}[\ell_L]=0$, this is exactly $\chi^2(D_L \| D_0)$.  The standard $\chi^2$-based sample-complexity bound (Le Cam's method) gives $m = \Omega(1/\chi^2) = \Omega(2^n)$ for decision.  For search, Fano's inequality with $|\Lagr| = 2^{n(n+1)/2+O(1)}$ and per-sample information $\Theta(2^{-n})$ yields $m = \Omega(n^2 \cdot 2^n)$.
\end{proof}

\paragraph{Consequence for reductions.} Propositions~\ref{prop:per-sample-mi} and~\ref{prop:chi2-sample} imply that \emph{membership-LSN with $\operatorname{poly}(n)$ samples is information-theoretically secure}: no reduction, linear or non-linear, adaptive or static, can recover $L$ or distinguish $D_L$ from $D_0$ with only polynomially many samples.  This makes ``membership-LSN $\not\to$ LPN with $\operatorname{poly}(n)$ samples'' trivially true and cryptographically empty.  The hardness assumption relevant to applications is \emph{batch-LSN with pseudorandom $A$} (\Cref{def:lsn-batch}).  The barrier studied by LPQR26 is \emph{sympLPN} (\Cref{def:symplpn}), a decoding problem with linear labels and public isotropic matrix.

\subsection{Distance Distribution and Expected Intersection}
\label{subsec:distance}

The correlation between two LSN distributions depends only on the dimension $j = \dim(L \cap L')$ of the intersection of their secrets. The distribution of $j$ is:

\begin{theorem}[Distance distribution]
\label{thm:distance}
For uniform $L, L' \in \Lagr(2n, \F_2)$,
\[
\Pr[j=k] = \frac{{n \brack k}_2 \cdot 2^{(n-k)(n-k+1)/2}}{|\Lagr(2n,\F_2)|}.
\]
Numerical evaluation shows that $\E[j] \to 0.76$ and $\operatorname{Var}(j) \to 0.60$ as $n \to \infty$.
\end{theorem}

The rapid decay of \Cref{eq:j-distribution} means that two random Lagrangians typically intersect in dimension $0$ or $1$; intersections of dimension $\geq 14$ occur with probability $< 2^{-100}$ for every $n \geq 14$ (the tail $\Pr[j=k] \approx 2^{-k(k+1)/2}$ is essentially independent of $n$).

\subsection{The Diluted-LPN View}
\label{subsec:diluted-lpn}

In the Siegel chart $L = \{(u, Mu) : u \in \F_2^n\}$ with $M \in \operatorname{Sym}(n,\F_2)$, a sample $(a,b)$ with $a=(u,v)$ and $b=1$ reveals $n$ linear equations $v = Mu$ in the $n(n+1)/2$ unknown entries of $M$.  The catch is that true positives are drowned in noise.

\begin{proposition}[Dilution of true positives]
\label{prop:dilution}
Conditioned on the label $b=1$, the posterior probability that the query point lies in $L$ is
\[
\Pr[a \in L \mid b=1] \;=\; \frac{(1-p)\,2^{-n}}{p + (1-2p)\,2^{-n}} \;=\; \Theta(2^{-n}).
\]
At $p=1/4$ the exact constant is $3$ asymptotically ($\approx 3\cdot 2^{-n}$ for large $n$); for $n=3$ it is exactly $3/10$.
\end{proposition}

\begin{proof}
Bayes' rule: $\Pr[a \in L \mid b=1] = \Pr[b=1 \mid a \in L]\Pr[a \in L] / \Pr[b=1] = (1-p)\cdot 2^{-n} / (p + (1-2p)2^{-n})$.
\end{proof}

\Cref{prop:dilution} reframes batch-LSN search as an \emph{$\varepsilon$-diluted LPN} problem over $\operatorname{Sym}(n,\F_2)$: with probability $\varepsilon \approx 3\cdot 2^{-n}$ a sample yields $n$ clean linear equations in $M$; otherwise it yields noise.  This is LPN with secret dimension $n(n+1)/2$ at an extreme dilution rate.  The computational question becomes:

\paragraph{Enrichment question.}
\label{q:enrichment}
Can any efficient selection rule (using past labeled samples and the symplectic structure) produce query points with $\Pr[a \in L] \geq (1+\delta)\,2^{-n}$ for non-negligible $\delta$?  In other words, can the $2^{-n}$ dilution be beaten?

A negative answer would be the cleanest quantitative statement of what $\mathrm{LSN}\setminus\mathrm{LPN}$ is; a positive answer would yield a concrete attack improvement (and a step toward a $6.5$th-family demotion).  The SQ lower bounds of the core paper suggest that enrichment is impossible for SQ-implementable selection rules: any such rule is a linear combination of SQ queries, and the correlation bounds limit its advantage to $2^{-n}$. The following lemma makes the reduction from enrichment to distinguishing explicit.

\begin{lemma}[Enrichment boosts SQ advantage]
\label{lem:enrichment-sq}
Let $\mathcal{R}$ be a selection rule that, using past samples from the membership-LSN oracle, outputs a query $a \in \F_2^{2n}$ such that $\Pr[a \in L] = (1+\delta)\,2^{-n}$ for some $\delta \ge 0$. Then the statistical query $\phi(a,b) = (-1)^b$ has advantage
\[
  \operatorname{Adv}(\phi) = 2|1-2p|(1+\delta)2^{-n}.
\]
In particular, the gain over the uniform-query baseline ($\delta=0$) is $\Delta\operatorname{Adv} = 2|1-2p|\cdot\delta\cdot 2^{-n}$.
\end{lemma}

\begin{proof}
Under the null model ($b=e$) we have $\E[\phi] = \E[(-1)^e] = 1-2p$. Under the structured model, $\Pr[b=1 \mid a] = p + \mathbf{1}_L(a)(1-2p)$, so $\E[\phi \mid a] = 1-2\Pr[b=1 \mid a] = (1-2p)(1-2\mathbf{1}_L(a))$. Averaging over $a \sim \mathcal{R}$ gives $\E[\phi] = (1-2p)(1-2\Pr[a \in L]) = (1-2p)(1-2(1+\delta)2^{-n})$. The advantage is $|\E[\phi]-\E_{\text{null}}[\phi]| = |1-2p| \cdot 2(1+\delta)2^{-n}$.
\end{proof}

A $\delta$-enrichment would scale the per-query advantage by $(1+\delta)$; for selection rules implementable by linear SQ this exceeds the per-query linear-SQ ceiling proved in the core paper, while general (non-linear, multi-sample) selection rules remain open.  Empirical probes confirm the predicted scaling $\delta \approx c\,m\,2^{-n}$ with $c=\Theta(1)$.  The exact constant depends on the prior, the sample-size regime, and the estimator; at $n=65$ and $m=22{,}528$ the enrichment is $\delta \le 2^{-50}$, far below any cryptographically relevant threshold.

\subsection{Parameters}
\label{subsec:parameters}

\Cref{tab:parameters} gives illustrative statistical-query-model levels derived from the conditional core-paper bound; these are not concrete-security estimates.

\begin{table}[ht]
\centering
\caption{Illustrative LSN statistical-query-exponent levels for $p=1/4$: the dimension $n$ at which the SQ-query lower bound reaches a target exponent $\lambda$ (i.e.\ $\log_2 q_{\min}\approx\lambda$). SQ-model query exponents, \emph{not} concrete bit-security levels.}
\label{tab:parameters}
\begin{tabular}{@{}cccc@{}}
\hline
Target exponent $\lambda$ & $n$ & $\log_2(q_{\min})^{\text{(a)}}$ & $|\Lagr|$ \\
\hline
$80$  & 41  & 80.4  & $\sim 2^{862}$  \\
$128$ & 65  & 128.4 & $\sim 2^{2146}$ \\
$192$ & 97  & 192.4 & $\sim 2^{4754}$ \\
$256$ & 129 & 256.4 & $\sim 2^{8386}$ \\
\hline
\end{tabular}
\vspace{4pt}
\noindent {\footnotesize $^{\text{(a)}}$ The $q_{\min}$ figures use the conditional SQ bound of the core paper; its unconditional spread bound gives $\Omega(2^n)$. These figures bound \emph{statistical-query} algorithms and are evidence, not a concrete-security guarantee.}
\end{table}

The $q_{\min}$ figures follow from the conditional bound $q \ge \tfrac13 \cdot 2^{2n-c}$ of the core paper (conditional on its pencil-extremality conjecture, whose absolute constant $c$ is normalized to $0$ for these illustrative levels), built on its exact pairwise-correlation formula: at $c=0$, $\log_2 q_{\min} = 2n - \log_2 3 \approx 2n - 1.585$, so the $\lambda=80$ line is met at $n=41$ with $0.4$ bits of slack, while an unknown $c>0$ would shift every level down uniformly. Since the multiplicative constant $\tfrac13$ is below $1$, parameters with $2n < \lambda + 1.6$ are not advised without a rigorous upper bound.

\section{Cryptographic Primitives}
\label{sec:primitives}

\subsection{LSN-SNARK}
\label{sec:snark}

We construct a preprocessing SNARK for the relation ``$x \in L$'' given a committed Lagrangian $L$, in the rank-1 constraint system (R1CS) / quadratic arithmetic program framework of Gennaro \etal\ \cite{GGPR13}, Parno \etal\ \cite{PHGR13}, Ben-Sasson \etal\ \cite{BCTV14}, and Groth \cite{Gro16}. The construction exploits the \emph{Siegel-chart} (graph-Lagrangian) representation to reduce the R1CS circuit from $O(n^3)$ to $O(n^2)$ constraints; here we count the R1CS constraints of the membership relation, which dominate proof and verification cost in these systems.

\paragraph{Siegel-chart representation.}
Fix a public symplectic matrix $S \in \mathrm{Sp}(2n,\F_2)$ (e.g., generated from a standard seed). Every ``generic'' Lagrangian can be written as
\[
L_M = S \cdot \{(x, Mx) : x \in \F_2^n\},
\]
where $M \in \mathrm{Sym}(n,\F_2)$ is a symmetric matrix. The set of graph Lagrangians has cardinality $|\mathrm{Sym}(n)| = 2^{n(n+1)/2}$, which differs from the full Lagrangian family $|\Lagr| = \prod_{i=1}^n (2^i+1)$ by only a constant factor ($\approx 2.38$). Consequently, restricting to graph Lagrangians reduces the key-space by only a constant factor, leaving the brute-force exponent $n^2/2+O(n)$ unchanged. We expect the SQ lower bounds of the core paper to carry over with the same exponent --- the pairwise-correlation structure is preserved under the symplectic change of coordinates $S$ --- although a separate worst-subset SDA bound for this graph sub-family is not proved here and is treated as evidence rather than theorem.

\paragraph{Binding $M$ to a public commitment.}
To prove knowledge of the secret $M$, the prover binds the witness to a public commitment $c = \mathsf{Hash}(M)$, where $\mathsf{Hash}$ is a SNARK-friendly hash function (e.g., Poseidon). This introduces preimage-resistance of $\mathsf{Hash}$ as a second assumption; we size $\mathsf{Hash}$ so that Grover preimage search dominates the LSN level --- a $2\lambda$-bit output gives $2^{\lambda}$ quantum preimage cost, so a $256$-bit Poseidon output keeps LSN the bottleneck at the $\lambda=128$ level, and an LSN instance with target exponent $\lambda$ needs a $\geq 2\lambda$-bit hash output in general. The SNARK circuit then verifies two things: (1)~commitment opening ($c = \mathsf{Hash}(M)$), and (2)~membership ($x \in L_M$). Hashing the $n(n+1)/2$ bits of $M$ with Poseidon adds only $O(n)$ constraints---packing the bits into $\approx 9$ field elements for $n=65$ costs roughly $300$--$500$ R1CS constraints. The full circuit size is therefore $n^2 + O(n)$, still $O(n^2)$ and still orders of magnitude below ZK-lattice circuits.

\paragraph{Membership circuit.}
The core membership sub-circuit verifies $x \in L_M$ in two steps:
\begin{enumerate}
\item Apply the public inverse $S^{-1}$ to $x$ to obtain $(a,b) = S^{-1}x$. This is a linear (XOR) operation over $\F_2$, hence \textbf{free} in a native-$\F_2$ system (over a prime field it adds only linear combinations, no multiplication gates).
\item Check $b = Ma$. Each entry $b_i = \sum_{j=1}^n M_{ij} a_j$ is a dot product of length $n$ over $\F_2$. Treating $a$ as a witness vector yields $n$ AND-gates per row, for $n^2$ \emph{native-$\F_2$ AND-gate} constraints total; over a prime-field R1CS each AND-gate is one multiplication constraint, with an additional $O(n^2)$ booleanity constraints for the witness bits. The final sums are linear and absorbed into the R1CS output side. Symmetry of $M$ is free: we encode only the upper-triangular $n(n+1)/2$ entries.

\paragraph{Further optimization.} If $a$ is provided directly as a public input (derived from $x$ by the verifier), each product $M_{ij}a_j$ becomes a selector---free when $a_j=0$ and a copy when $a_j=1$. The check $b_i = \sum_j M_{ij}a_j$ then collapses to a single linear constraint per row, giving \textbf{$n$ constraints} over $\F_2$ R1CS, or $n(n+3)/2$ constraints over a prime field with boolean witnesses for $M$.
\end{enumerate}

\Cref{tab:r1cs} compares LSN-SNARK with other ZK-PQC approaches. The constraint count is the dominant cost metric across preprocessing SNARKs \cite{GGPR13,Gro16} and transparent systems such as Bulletproofs \cite{BBB+18} alike.

\begin{table}[ht]
\centering
\caption{R1CS constraints for LSN membership vs.~ZK-PQC full primitives.}
\label{tab:r1cs}
{\small\begin{tabular}{@{}lcc@{}}
\hline
Scheme & Constraints ($n=65$) & Asymptotic \\
\hline
\textbf{LSN-SNARK membership} & $\mathbf{\approx 4{,}225}$ & $\mathbf{O(n^2)}$ \\
\textbf{LSN-SNARK full ($+\,$hash)} & $\mathbf{\approx 4{,}500}$--$\mathbf{4{,}700}$ & $\mathbf{O(n^2)}$ \\
SHA-256 (per block) & $\approx 25{,}000$ & $O(1)$ \\
AES-128 & $\approx 16{,}000$ & $O(1)$ \\
ZK-lattice KEM (full verif.) & millions & $O(n \log n)$ \\
ZK-SPHINCS+ (full sig. verif.) & millions & $O(\lambda \log \lambda)$ \\
\hline
\end{tabular}}
\end{table}

\paragraph{Witness structure.} The witness $\mathbf{w}$ consists of:
\begin{itemize}
\item $n(n+1)/2$ bits for the upper-triangular entries of $M \in \mathrm{Sym}(n,\F_2)$;
\item $n$ bits for the vector $a$ (the first half of $S^{-1}x$, computed as a linear function of $x$ but included in the witness for the dot-product checks);
\item auxiliary opening data for the hash commitment (omitted for deterministic hashes, negligible otherwise).
\end{itemize}
Total witness size: $n(n+1)/2 + O(n) = O(n^2)$ bits.

\paragraph{Verified counts.}
The $n^2$ prediction is exact for the membership sub-circuit when $a$ is a witness variable: $n=8$ yields $64$ constraints, $n=41$ ($\lambda{=}80$) yields $1{,}681$, and $n=65$ ($\lambda{=}128$) yields $4{,}225$. With public $a$, the optimized membership counts are $8$, $41$, and $65$ respectively over $\F_2$. The \emph{full} circuit (membership $+$ Poseidon commitment opening) at $n=65$ is $\approx 4{,}500$--$4{,}700$ constraints depending on hash parameterization.

\paragraph{Why this matters.} The membership relation $x \in L_M$ thus has an unusually compact ZK circuit---$\approx 4{,}225$ R1CS constraints at $n=65$ (or $65$ with public $a$), small relative to hashing one SHA-256 block ($\approx 25{,}000$). The full primitive (membership $+$ commitment opening) stays below $\approx 5{,}000$ constraints---about $5\times$ smaller than a single SHA-256 block ($\approx 25{,}000$) and orders of magnitude smaller than published ZK-lattice KEM verification circuits (reported in the millions of constraints; we do not benchmark a specific implementation). \textbf{Scope.} The table compares \emph{core-relation} circuits; full signature verification includes additional binding overhead that is comparable across schemes. The LSN counts are native-$\F_2$ AND-gates, whereas the lattice/hash comparisons are prime-field constraints; normalizing the $\F_2$ counts to a prime field adds $O(n^2)$ booleanity constraints, which does not change the order-of-magnitude gap. Our claim is not a record for the full signature, but that the \emph{membership} sub-circuit---the inner relation being proved---is exceptionally compact. That compactness helps building blocks such as \emph{verifiable PQ encryption} and \emph{recursive SNARK aggregation}, where many membership checks dominate the cost.

\paragraph{Batch verification.}
The compactness enables dramatic amortization when proving membership for $k$ distinct elements $x^{(1)},\dots,x^{(k)}$ under the same secret $M$. The commitment $c=\mathsf{Hash}(M)$ is verified \textbf{once}; each additional membership check adds only $n^2$ constraints (or $n$ with public $a$). For batch size $k$:
\[
\text{Total constraints} = \underbrace{O(n)}_{\text{hash opening}} + k \cdot \underbrace{O(n^2)}_{\text{membership}}.
\]
By contrast, ZK verification of lattice KEMs involves per-proof NTT-style evaluations and is not known to amortize comparably. Concrete example ($n=65$, $k=1000$): LSN uses $\approx 500 + 1000\cdot 4{,}225 \approx 4.2\times 10^6$ constraints; published lattice-KEM verification circuits are reported in the millions of constraints per proof, so a like-for-like batch would be substantially larger, but the precise ratio depends on the proving system, arithmetization, and the exact relation, which we do not benchmark here.

\subsection{On key encapsulation}
\label{sec:kem}

A natural goal is a key-encapsulation mechanism (KEM) from LSN. We deliberately do \emph{not} include one. The natural reconciliation construction we explored masks the encapsulated codeword with a value computable from public data, so it offers no security; and the constant-noise regime ($p=\tfrac14$) blocks the standard sparse-secret dual-LPN route, since the secret weight forced by syndrome-decoding hardness drives the reconciliation channel to capacity zero. A secure LSN-KEM therefore appears to require either dedicated constant-noise LPN-PKE machinery~\cite{YZ16} or a low-noise LSN variant; we leave a sound construction and its proof to future work. This affects neither the hardness foundation nor the SNARK above.

\begin{thebibliography}{99}

\bibitem[AAC+22]{AAC+22}
C.~Aguilar~Melchor, N.~Aragon, S.~Bettaieb, L.~Bidoux, O.~Blazy, J.-C.~Deneuville, P.~Gaborit, S.~Ghosh, A.~Hauteville, and G.~Z\'emor.
HQC.
NIST PQC Round 4 Submission, 2022.

\bibitem[ABD+21]{ABD+21}
M.~Albrecht, D.~Bernstein, T.~Chou, C.~Cid, J.~Gilcher, T.~Lange, V.~Maram, I.~von Maurich, R.~Misoczki, R.~Niederhagen, K.~Paterson, E.~Persichetti, C.~Peters, P.~Schwabe, N.~Sendrier, J.~Szefer, C.~Tjhai, M.~Tomlinson, and W.~Wang.
Classic McEliece.
NIST PQC Round 4 Submission, 2022.

\bibitem[ABW10]{ABW10}
B.~Applebaum, B.~Barak, and A.~Wigderson.
Public-key cryptography from different assumptions.
In \emph{STOC}, pages 171--180. ACM, 2010.

\bibitem[Ale11]{Ale11}
M.~Alekhnovich.
More on average case vs approximation complexity.
\emph{Computational Complexity}, 20(4):755--786, 2011.

\bibitem[BBB+18]{BBB+18}
B.~B\"unz, J.~Bootle, D.~Boneh, A.~Poelstra, P.~Wuille, and G.~Maxwell.
Bulletproofs: Short proofs for confidential transactions and more.
In \emph{IEEE S\&P}, pages 315--334, 2018.

\bibitem[BCTV14]{BCTV14}
E.~Ben-Sasson, A.~Chiesa, E.~Tromer, and M.~Virza.
Scalable zero knowledge via cycles of elliptic curves.
In \emph{CRYPTO}, pages 276--294. Springer, 2014.

\bibitem[BFJ+94]{BFJ+94}
A.~Blum, M.~Furst, J.~Jackson, M.~Kearns, Y.~Mansour, and S.~Rudich.
Weakly learning DNF and characterizing statistical query learning using Fourier analysis.
In \emph{STOC}, pages 253--262. ACM, 1994.

\bibitem[BFKL94]{BFKL94}
A.~Blum, M.~Furst, M.~Kearns, and R.~Lipton.
Cryptographic primitives based on hard learning problems.
In \emph{CRYPTO}, pages 278--291. Springer, 1994.

\bibitem[BKW03]{BKW03}
A.~Blum, A.~Kalai, and H.~Wasserman.
Noise-tolerant learning, the parity problem, and the statistical query model.
\emph{J.\ ACM}, 50(4):506--519, 2003.

\bibitem[BPR12]{BPR12}
A.~Banerjee, C.~Peikert, and A.~Rosen.
Pseudorandom functions and lattices.
In \emph{EUROCRYPT}, pages 719--737. Springer, 2012.

\bibitem[BCGI18]{BCGI18}
E.~Boyle, G.~Couteau, N.~Gilboa, and Y.~Ishai.
Compressing vector OLE.
In \emph{ACM CCS}, pages 896--912, 2018.

\bibitem[CRSS98]{CRSS98}
A.~R.~Calderbank, E.~M.~Rains, P.~W.~Shor, and N.~J.~A.~Sloane.
Quantum error correction via codes over GF(4).
\emph{IEEE Trans.\ Inform.\ Theory}, 44(4):1369--1387, 1998.

\bibitem[DKS17]{DKS17}
I.~Diakonikolas, D.~M.~Kane, and A.~Stewart.
Statistical query lower bounds for robust estimation of high-dimensional Gaussians and Gaussian mixtures.
In \emph{FOCS}, pages 73--84. IEEE, 2017.

\bibitem[DGM12]{DGM12}
N.~D\"ottling, J.~M\"uller-Quade, and A.~Nascimento.
IND-CCA secure cryptography based on a variant of the LPN problem.
In \emph{ASIACRYPT}, pages 485--503. Springer, 2012.

\bibitem[EKM17]{EKM17}
A.~Esser, R.~K\"ubler, and A.~May.
LPN decoded.
In \emph{CRYPTO}, pages 486--514. Springer, 2017.

\bibitem[FGR+17]{FGR+17}
V.~Feldman, E.~Grigorescu, L.~Reyzin, S.~S.~Vempala, and Y.~Xiao.
Statistical algorithms and a lower bound for detecting planted cliques.
\emph{Journal of the ACM}, 64(2):8:1--8:37, 2017.

\bibitem[GGPR13]{GGPR13}
R.~Gennaro, C.~Gentry, B.~Parno, and M.~Raykova.
Quadratic span programs and succinct NIZKs without PCPs.
In \emph{EUROCRYPT}, pages 626--645. Springer, 2013.

\bibitem[GJL14]{GJL14}
Q.~Guo, T.~Johansson, and C.~L\"ondahl.
Solving LPN using covering codes.
In \emph{ASIACRYPT}, pages 1--20. Springer, 2014.

\bibitem[GL22]{GL22}
A.~Gollakota and D.~Liang.
On the hardness of PAC-learning stabilizer states with noise.
\emph{Quantum}, 6:640, 2022. arXiv:2102.05174.

\bibitem[Got97]{Got97}
D.~Gottesman.
Stabilizer codes and quantum error correction.
PhD thesis, Caltech, 1997.

\bibitem[Gro96]{Gro96}
L.~K.~Grover.
A fast quantum mechanical algorithm for database search.
In \emph{STOC}, pages 212--219. ACM, 1996.

\bibitem[Gro16]{Gro16}
J.~Groth.
On the size of pairing-based non-interactive arguments.
In \emph{EUROCRYPT}, pages 305--326. Springer, 2016.

\bibitem[HKL+12]{HKL+12}
S.~Heyse, E.~Kiltz, V.~Lyubashevsky, C.~Paar, and K.~Pietrzak.
Lapin: An efficient authentication protocol based on ring-LPN.
In \emph{FSE}, pages 346--365. Springer, 2012.

\bibitem[Kea98]{Kea98}
M.~Kearns.
Efficient noise-tolerant learning from statistical queries.
\emph{J.\ ACM}, 45(6):983--1006, 1998.

\bibitem[Kir11]{Kir11}
P.~Kirchner.
Improved generalized birthday attack.
\emph{Cryptology ePrint Archive}, Report 2011/377, 2011.

\bibitem[KLP+25]{KLP+25}
A.~Khesin, J.~Lu, A.~Poremba, A.~Ramkumar, and V.~Vaikuntanathan.
Average-Case Complexity of Quantum Stabilizer Decoding.
\emph{arXiv:2509.20697}, 2025.

\bibitem[LF06]{LF06}
E.~Levieil and P.-A.~Fouque.
An improved LPN algorithm.
In \emph{SCN}, pages 348--359. Springer, 2006.

\bibitem[LPQR26]{LPQR26}
J.~Lu, A.~Poremba, Y.~Quek, and A.~Ramkumar.
Post-quantum cryptography from quantum stabilizer decoding.
\emph{arXiv:2603.19110}, 2026.

\bibitem[Lyu05]{Lyu05}
V.~Lyubashevsky.
The parity problem in the presence of noise, decoding random linear codes, and the subset sum problem.
In \emph{APPROX-RANDOM}, pages 378--389. Springer, 2005.

\bibitem[AG04]{AG04}
S.~Aaronson and D.~Gottesman.
Improved simulation of stabilizer circuits.
\emph{Phys.\ Rev.\ A}, 70:052328, 2004. arXiv:quant-ph/0406196.

\bibitem[McE78]{McE78}
R.~J.~McEliece.
A public-key cryptosystem based on algebraic coding theory.
\emph{DSN Progress Report 42--44}, pages 114--116, 1978.

\bibitem[NC00]{NC00}
M.~A.~Nielsen and I.~L.~Chuang.
Quantum Computation and Quantum Information.
Cambridge University Press, 2000.

\bibitem[NIS22]{NIS22}
NIST.
Post-quantum cryptography standardization, 2022.
\url{https://csrc.nist.gov/projects/post-quantum-cryptography}.

\bibitem[PHGR13]{PHGR13}
B.~Parno, J.~Howell, C.~Gentry, and M.~Raykova.
Pinocchio: Nearly practical verifiable computation.
In \emph{IEEE S\&P}, pages 238--252, 2013.

\bibitem[PQS26]{PQS26}
A.~Poremba, Y.~Quek, and P.~Shor.
The Learning Stabilizers with Noise problem.
In \emph{ITCS}, LIPIcs vol.~362, pages 108:1--108:19, 2026. arXiv:2410.18953.

\bibitem[Pra62]{Pra62}
E.~Prange.
The use of information sets in decoding cyclic codes.
\emph{IRE Trans.\ Inform.\ Theory}, 8(5):5--9, 1962.

\bibitem[Reg05]{Reg05}
O.~Regev.
On lattices, learning with errors, random linear codes, and cryptography.
\emph{J.\ ACM}, 56(6):34, 2009.

\bibitem[SAB+22]{SAB+22}
P.~Schwabe, R.~Avanzi, J.~Bos, L.~Ducas, E.~Kiltz, T.~Lepoint, V.~Lyubashevsky, J.~M.~Schanck, G.~Seiler, and D.~Stehl\'e.
CRYSTALS-Kyber.
NIST PQC Round 3 Submission, 2022.

\bibitem[Sho94]{Sho94}
P.~W.~Shor.
Algorithms for quantum computation: Discrete logarithms and factoring.
In \emph{FOCS}, pages 124--134. IEEE, 1994.

\bibitem[YZ16]{YZ16}
Y.~Yu and J.~Zhang.
Cryptography with auxiliary input and trapdoor from constant-noise LPN.
In \emph{CRYPTO}. Springer, 2016. IACR ePrint 2016/514.

\end{thebibliography}

\end{document}
